\section{品質を工程と判断の関係で考える}
\label{theory:quality:background}
品質管理で問うのは、ばらつきの小ささだけではない。何を満たすべきかという要求、
どのように作るかという工程、どの情報で異常を判断するかという観測を結ぶことである。
良い製品を作ること、工程が安定していること、出荷ロットを受け入れることは判断対象が違う。
目標を設定し、工程を観測して改善する流れの中で、確率統計は限られた測定から判断する道具となる。

\section{平均・分散・共分散から合成する}
\label{theory:quality:variance}
有限な二次モーメントを持つ変数を列ベクトル \(X\)、係数を \(a\) とする。
線形結合 \(Y=a^{\mathsf T}X\) の平均は \(E[Y]=a^{\mathsf T}E[X]\) である。
平均からの偏差を二乗して展開すると
\[
 \operatorname{Var}(Y)
 =\sum_i a_i^2\sigma_i^2+2\sum_{i<j}a_i a_j\operatorname{Cov}(X_i,X_j)
 =a^{\mathsf T}\Sigma a
\]
となる。\(\sigma_i^2\) は成分の分散、\(\Sigma\) は共分散行列である。
平均の線形性には独立性は不要である。分散から共分散項を消すには無相関が必要で、
独立性はその十分条件となる。独立な差 \(X-Y\) でも係数 \(-1\) が二乗されるため
分散は和となる。共通加工機の温度変動などがあれば独立とは限らず、共分散を捨てない。
単位も二乗され、長さの分散は長さの二乗、標準偏差は長さである。\cite{stat_psu_linear}

\section{正規分布の閉性と上側確率}
\label{theory:quality:normal}
\(X\sim N(\mu,\sigma^2)\) の第二引数は分散とする。
独立な正規変数の線形結合も正規となる。これは特性関数
\(\exp(i\mu t-\sigma^2t^2/2)\) の積を取れば、平均は係数付き和、
分散は係数の二乗付き和になることから分かる。
一般には各成分が正規であるだけでは足りず、同時分布の条件も必要である。
標準化 \(Z=(X-\mu)/\sigma\) により
\[
 \Pr(X\ge x)=1-\Phi\!\left(\frac{x-\mu}{\sigma}\right)
\]
と表せる。\(\Phi\) は標準正規累積分布であり、右側の面積を求める。
試験表が「上側確率」を示すか、左側の累積確率を示すかを先に読む。
差の平均に符号を付け、分散を合成し、平方根で標準偏差を出してから標準化する。\cite{stat_nist_normal}

\section{標本平均と総量は別の変数である}
\label{theory:quality:sample_mean}
同じ平均 \(\mu\)、分散 \(\sigma^2\) を持つ独立な \(n\) 個について、
総量の期待値は \(n\mu\)、分散は \(n\sigma^2\) である。
標本平均は総量を \(n\) で割るため、期待値は \(\mu\)、分散は
\(\sigma^2/n\)、標準誤差は \(\sigma/\sqrt n\) となる。
物理量の平均を取っても期待値を標本数で割って小さくするわけではない。
原変数が正規なら有限の \(n\) でもこの平均は正規である。
中心極限定理による大標本の近似とは別の根拠である。
相関のある標本では共分散項が残り、標本数を増やすだけで同じ割合で不確かさが減るとは限らない。

\section{応力・強度から故障確率へ}
\label{theory:quality:stress_strength}
強度 \(R\) と作用応力 \(S\) の単位を揃え、限界状態 \(G=R-S\) を定める。
独立正規モデルでは
\[
 b=\frac{\mu_R-\mu_S}{\sqrt{\sigma_R^2+\sigma_S^2}},\qquad
 \Pr(G<0)=\Phi(-b)
\]
となる。\(b\) は平均の余裕をばらつきで規格化した無次元の指標である。
平均強度が大きくてもばらつきが大きければ故障確率は大きくなり得る。
安全率の比とこの確率を同一視しない。
非正規の裾や強度と応力の依存があると、この式の直接適用はできない。
接続部のすき間も内径から外径を引く限界状態であり、平均の大小は全個体の符号を保証しない。

\section{どの量を数えるかで分布を選ぶ}
\label{theory:quality:models}
同じ「ばらつき」でも、測る量の範囲が違う。公平なサイコロは有限集合の離散一様である。
一定の成功確率 \(p\) の独立試行を \(n\) 回行った成功数は二項分布で、
\[
 \Pr(K=k)={n\choose k}p^k(1-p)^{n-k},\qquad
 E[K]=np,\quad \operatorname{Var}(K)=np(1-p)
\]
である。製造数を固定して不適合品を数えるモデルもこれに対応する。
コインが公平でなくても、一定 \(p\) と独立性があれば二項である。
品質の共通原因による連発や、試行ごとの確率変化があればこの条件を満たさない。
\cite{stat_nist_binomial}

\section{ポアソンの回数と指数の待ち時間}
\label{theory:quality:poisson}
二項分布で機会を増やし、小確率にして \(np=\lambda\) を有限に保つと、
個数は \(\Pr(K=k)=e^{-\lambda}\lambda^k/k!\) のポアソンへ近づく。
平均と分散はともに \(\lambda\) である。
一定率 \(r\) のポアソン過程なら時間 \(t\) の回数の平均は \(rt\) で、
一度も生じない確率が \(e^{-rt}\) となる。
したがって次の事象までの待ち時間 \(T\) は
\(\Pr(T>t)=e^{-rt}\) の指数分布となり、平均は \(1/r\) である。
率 \(r\) と区間内平均 \(\lambda=rt\) の単位を区別する。
災害や交通事故の実測に対する適合は別に検証する必要がある。
大きな \(\lambda\) ではポアソンの歪度 \(1/\sqrt{\lambda}\) が小さくなり、
平均 \(\lambda\)、分散 \(\lambda\) の正規近似が有効になり得る。
「回数だから正規近似は常に不可能」とはしない。\cite{stat_nist_poisson,stat_nist_exponential}

\section{要求品質を設計特性へ展開する}
\label{theory:quality:qfd}
利用者の言葉である要求品質を、設計で調整する品質要素へ変換するのが品質機能展開である。
試験の品質表では要求 \(i\) の重要度 \(w_i\) と要素 \(j\) との関係係数 \(c_{ij}\) を使い、
\(v_j=\sum_i w_i c_{ij}\) を求める。単なる丸の数ではなく、重い要求との対応が寄与する。
係数尺度は判断の結果に影響するため、指定された◎・○・△の値を守る。
この加重得点は実測の成功確率や安全性そのものではない。
関連する特性間の相反もあり、得点最大の列だけを設計の唯一の目標にはしない。
\cite{quality_qfd_scale}

\section{検査と上流の品質保証}
\label{theory:quality:assurance}
検査は要求への適合を判定する活動である。検査で不適合を見つけても、
それだけで不適合を作る原因は除かれない。源流管理は仕様や工程など上流にさかのぼって
発生原因を抑える考え方である。
標準化で手順と判断を明確にし、維持管理と改善を結ぶ。
工程の標準は経験的な固定ルールではなく、要求・観測・変更理由を追跡できるものとする。
QC工程表は工程と管理項目などを結ぶ計画の道具であり、統計量を時間順に並べる管理図と違う。
\cite{quality_process_jsqc}

\section{PDCAをフィードバックとして読む}
\label{theory:quality:pdca}
Planは目標と達成するプロセスを定め、Doは実施、Checkは方針・目標・要求に対する
結果を監視・測定し、Actは性能を改善する行動を取る段階である。
目標なしの測定、測定なしの改善では循環の根拠が失われる。
CをChangeやChoose、AをAdjustやAchieveという近い一般語へ置き換えない。
Actの結果を次の計画と標準へ戻すことに技術的意味がある。
\cite{quality_iso_process}

\section{2008年版を扱う問題の資料時点}
\label{theory:quality:historical}
H27の設問はISO 9001:2008/JIS Q 9001:2008を指定する。
ISOの当時の実施ガイダンスN836でもPlan-Do-Check-Actの利用が示される。
後版である2015年版の公式解説は、プロセスの管理とリスクに基づく思考を併せて説明する。
ここでは一般の四段階を両資料で照合したのであり、2008年版の全条文を取得・照合した
とは扱わない。現行規格の詳細を求める場合はその版・追補を別途確認する。
\cite{quality_iso2008_guidance,quality_iso_process}

\section{工程能力と規格限界}
\label{theory:quality:capability}
工程能力は、安定した工程が作る特性のばらつきと規格の許容幅を比較する。
正規モデルで平均 \(\mu\)、標準偏差 \(\sigma\)、規格上限 \(U\)、下限 \(L\) とすると
\[
 C_p=\frac{U-L}{6\sigma},\qquad
 C_{pk}=\min\left(\frac{U-\mu}{3\sigma},\frac{\mu-L}{3\sigma}\right)
\]
である。\(C_p\) は幅の比較だけなので平均の片寄りを直接反映しない。
\(C_{pk}\) は近い規格限界までの余裕を見る。
両者は無次元であり、標準偏差そのものではない。
非正規、工程の漂移、異なる集団の混合では、指数から不適合率への読み替えに注意する。
一定の数値を全産業共通の合格基準として勝手に設定しない。\cite{quality_nist_capability}

\section{管理図・ヒストグラム・散布図}
\label{theory:quality:charts}
管理図は工程の特性又は統計量を時間・標本順に並べ、安定工程の変動を基準に異常を探す。
中心線と管理限界は観測モデルに基づき、利用者の規格限界とは目的が違う。
限界内に全点があっても、片側への連続や系統的変化があれば安定とは限らない。
ヒストグラムは一変数の測定値を区間に分けて分布を示す。
区間幅が異なるなら高さを度数密度にし、長方形の面積を度数に対応させる。
散布図は対になった二変数の関連を示す。
分布、時間変動、二変数の関連という軸を分ければ、図の名を暗記するだけでなく用途が分かる。
\cite{quality_nist_control,quality_nist_histogram}

\section{相関は線形関連を測る}
\label{theory:quality:correlation}
有限かつ非零の分散について相関は
\(\rho=\operatorname{Cov}(X,Y)/(\sigma_X\sigma_Y)\) であり、
Cauchy--Schwarzから \(-1\le\rho\le1\) となる。
端点は完全な線形関係で達成されるため除外しない。
独立なら無相関だが、無相関だけで独立とは限らない。
対称な \(X\) と \(Y=X^2\) はその反例で、三次モーメントが零でも関数の依存が残る。
相関の絶対値は線形関連の強さを表し、因果の方向や共通原因を判別しない。
正の相関が1へ近づく説明を、任意のデータが完全な直線上にあるという断定へ変えない。

\section{回帰と決定係数の条件}
\label{theory:quality:regression}
回帰は説明変数から応答を予測するモデルを作り、残差などで適合を評価する。
切片を含む一説明変数の通常の最小二乗法では、偏差平方和を
回帰が説明する部分と残差の部分に分けられ、決定係数は \(R^2=r^2\) である。
相関の絶対値 \(|r|\) とは違う。
切片なし、多変数、異なるモデルで同じ式を機械的に使わない。
主成分分析は複数変数を合成して分散を捉える次元縮約であり、指定された応答を予測する
回帰と目的が違う。相関と回帰のどちらも観測だけから因果の証明にはならない。
\cite{stat_psu_correlation}

\section{ロットと抜取検査の二つの危険}
\label{theory:quality:acceptance}
全数を検査できない場合、サンプルからロットを受け入れるか判断する。
計数値では不適合品などを数え、計量値では長さなどの連続量を使う。
良いロットを拒否する生産者危険と、悪いロットを受け入れる消費者危険を分ける。
母集団の不適合率 \(p\) を決めた条件付き受入確率を扱うので、
「受け入れたロットが良い確率」という逆向きの条件付き確率とは違う。
\cite{quality_nist_sampling}

\section{OC曲線から閾値の効果を導く}
\label{theory:quality:oc}
大きなロットから近似的に独立に \(n\) 個を抽出し、不適合数 \(K\le c\) なら合格とする。
受入確率は二項の累積で
\[
 P_a(p)=\sum_{k=0}^c {n\choose k}p^k(1-p)^{n-k}
\]
である。良い品質点を \(p_{\mathrm g}\)、悪い品質点を \(p_{\mathrm b}\) とすれば、
\(\alpha=1-P_a(p_{\mathrm g})\)、\(\beta=P_a(p_{\mathrm b})\) となる。
\(n\) を固定して \(c\) を増やすと受入集合が広がるため、\(\alpha\) は下がり
\(\beta\) は上がる。この包含関係がトレードオフの根拠である。
\((n,c)\) の両方を固定すれば危険も固定であり、原典の「手順固定」の範囲は注意を要する。
小さい有限ロットからの非復元抽出は超幾何となり、二項の独立近似を無条件に使わない。
\cite{quality_nist_oc}

\section{ロット間の切換えと複数回の抽出}
\label{theory:quality:switching}
先行ロットの成績で後続の検査をなみ検査・きつい検査へ切り換える考え方と、
一つのロットで段階的に標本を追加する多回抜取を区別する。
後者は標本を追加したうえで合否を決める手順である。
切換えの回数や復帰条件は適用する規格に従うものであり、今回取得していない
出題時のJIS全文の数値規則を補わない。\cite{quality_nist_sampling}

\section{仮説検定は誤りを管理する判断である}
\label{theory:quality:hypothesis}
帰無仮説 \(H_0\) を基準とし、対立仮説に対してどれほど不整合な観測かを判断する。
真の \(H_0\) を棄却するのが第1種、偽の \(H_0\) を棄却しないのが第2種の誤りである。
有意水準 \(\alpha\) は第1種の誤りを管理する事前の基準であり、
観測後に \(H_0\) が真である確率ではない。
棄却されなかったことも \(H_0\) の真の証明にならない。
第2種の確率 \(\beta\) は具体的な対立側の値に依存し、検出力は \(1-\beta\) である。
同じ二種類の誤りが、品質の受入判断では生産者と消費者の危険に対応する。
\cite{quality_nist_sampling}

\section{t統計量と信頼区間の被覆}
\label{theory:quality:t_interval}
独立な正規標本で母分散が未知なら、不偏分散
\(s^2=\sum_i(X_i-\bar X)^2/(n-1)\) を使う。
平均を推定した残差は和が零という一つの制約を持つため自由度は \(n-1\) となる。
\[
 T=\frac{\bar X-\mu_0}{s/\sqrt n},\qquad
 \bar X\pm t_{1-\alpha/2,n-1}\frac{s}{\sqrt n}
\]
は、順に \(H_0\) 下での一標本t統計量と母平均の両側信頼区間である。
信頼係数は \(1-\alpha\) であって \(\alpha\) の逆数ではない。
反復標本で作る区間のうち母平均を覆う割合を表すので、観測後の一つの固定区間について
母平均がそこに存在する事後確率とはしない。\cite{stat_nist_interval}

\section{典型的な問われ方}
寸法・力の計算は変数の和差と分散の合成を先に決め、正規表で候補の範囲を判定する。
品質表は指定された関係尺度を行の重要度と掛け、列を合計する。
知識問題は要求適合、工程監視、工程能力、受入検査の判断対象を分ける。
各年度の組合せ表の列と番号は原典で照合し、同じ理論でも候補番号が同じと仮定しない。
ここで扱う例は今回確認した問題群であり、全試験回の頻度の確定は別の集計に従う。
